% 原书第 27—31 页；章首导言及第 3.1 节。
\chapter{序贯超级学习}
\label{ch:sequential-super-learning}
\chapterauthors{Sherri Rose 与 Mark J. van der Laan}
\chapteraffiliations{%
S. Rose（通讯作者）：美国哈佛医学院卫生保健政策系，180 Longwood Ave, Boston, MA 02115, USA；电子邮箱：\href{mailto:rose@hcp.med.harvard.edu}{rose@hcp.med.harvard.edu}。\\
M. J. van der Laan：美国加利福尼亚大学伯克利分校生物统计学部及统计学系，101 Haviland Hall, \#7358, Berkeley, CA 94720, USA；电子邮箱：\href{mailto:laan@berkeley.edu}{laan@berkeley.edu}。\\
本章 DOI：\url{https://doi.org/10.1007/978-3-319-65304-4_3}。}

假设一位医生希望预测：一组患者在两种特定用药方案下，未来两个时点的个体结局。她所问的也就是：如果我让这些患者接受用药方案 1 或用药方案 2，他们每个人在未来这两个时点的结局会怎样？哪种处理能让患者获得更好的疗效结局？哪种处理能让患者获得更好的安全性结局？一般而言，预测问题本身就可能具有纵向性质；我们常常希望了解，具有某些特征的患者在数月或数年后的平均结局会是什么样。我们通常是在这样一种假想情境下提出这个问题的：为患者指定某个特定的处理“规则”。

本章讨论纵向数据（longitudinal data）问题中的序贯超级学习（sequential super learning），为任意时点的预测提供一个框架。在《靶向学习》（\emph{Targeted Learning}，2011）第 3 章中，我们介绍了简单数据结构 $O=(W,A,Y)\sim P_0$ 下用于预测的超级学习（super learning），并估计条件均值 $E[Y\mid A,W]$。现在，我们要研究纵向数据结构和更复杂的预测问题。这类纵向数据问题，与我们的 $E[Y\mid A,W]$ 预测问题有一个自然的对应：$E(Y^d\mid W)=E(Y\mid A=d(W),W)$。在 $O=(W,A,Y)$ 中，基线之后只有一个观测结局，也只有一个干预节点，因此预测问题很简单：只需进行一次回归（以超级学习来估计）。这里，我们引入序贯预测（sequential prediction）的概念，将《靶向学习》（2011）第 3 章介绍的概念和方法，推广到多个时点实施动态干预 $d$ 的情形，以估计给定基线协变量时，反事实结局 $Y^d$ 的条件均值。

对于这些多时点干预，可以用序贯回归（sequential regression）来估计 $E(Y^d\mid L_0)$；这些序贯回归，也是利用纵向靶向最小损失估计（longitudinal targeted minimum loss-based estimation，LTMLE）估计 $E(Y^d)$ 及其他参数时的关键组成部分。我们可以直接对 $E(Y^d\mid L_0)$ 的估计值取平均，以估计 $E(Y^d)$。不过，下一章将讨论如何在 LTMLE 的每一步中实施靶向更新，改进 $E(Y^d)$ 的序贯回归估计量。

因此，我们要记住，预测问题往往就是科学研究的目标，序贯超级学习是适用于这类问题的独立工具。同时，我们也强调，所关注的还包括估计数据概率分布的目标参数（target parameter）；它通常可以解释为因果效应。在这类情境下，我们将实施 LTMLE。对于涉及纵向数据的研究问题，这一估计过程不可或缺的一环，就是对 $P_0$ 中计算目标参数所需的相关部分 $Q$ 进行序贯回归估计。本章通过序贯超级学习介绍这一步骤，第 4 章则介绍 LTMLE。因此，本章聚焦于靶向学习路线图中的条件均值估计，它对预测问题和因果效应问题都很有用。

\begin{sourcebox}[纵向数据中的效应估计与预测]
纵向数据中的因果效应研究问题和预测研究问题，本质上都是估计问题。对于前者，我们关注的是估计某项干预或某个动态过程的因果效应，或其他纵向效应问题。对于预测，我们希望构造一个函数：输入变量 $(A_{0:t},L_{0:t})$，即可预测结局的取值，其中也可能涉及动态过程。尽管两者常被混为一谈，但它们显然是不同的研究问题。LTMLE 的过程中包含预测步骤；因此，理解用于纵向数据预测的序贯超级学习器，是研究这两类问题都必须掌握的核心概念。
\end{sourcebox}

\section{背景：集成学习}
\label{sec:ch03-ensemble-background}

如第 1 章所述，集成学习（ensemble learning）已经针对多种数据类型和研究问题得到发展，并已应用于许多实质性研究领域。其核心预测框架是一种“集成式”的超级（机器）学习方法：它借助交叉验证（cross-validation），得到多个算法的最优加权平均。这一方法解决了预测中的关键难题：从决策树、惩罚回归到神经网络，可用的算法有很多，而在一个给定的数据集中，各个算法的表现可能很不一样。

长期以来的问题是：我们如何事先知道，哪种算法在我们的数据上表现最好？即使数据类型相似，也可能出现这样的情况：在一项研究中，设定错误的参数模型中的 logistic 回归优于决策树，而在另一项研究中，决策树又优于 logistic 回归。数据科学文献中有许多这样的例子。

既然采用的是非参数统计模型（nonparametric statistical model），我们起初可能会倾向于使用在不过拟合的前提下对数据进行平滑的非参数方法。然而，即使是局部平均这样简单的非参数估计量，也仍然需要对协变量空间进行划分，才能确定这一回归估计量的平滑程度。即便选择了最优划分，如果真实的底层数据生成分布（data-generating distribution）十分光滑，设定错误的参数统计模型中的 logistic 回归，也可能优于这个非参数估计量。

这些考虑促成了大量关于集成的统计学研究，并最终催生了超级学习器。超级学习避免了事先选定某个表现不佳的单一算法的风险。我们转而考虑许多算法，同时纳入局部平均、logistic 回归以及其他众多算法，也就不必担心局部平均会被 logistic 回归击败。这是因为，超级学习依据一个预先指定的损失函数（loss function），构造出所有候选算法的最优加权平均预测函数。

\medskip
\noindent\textbf{有代表性的应用}
\begin{itemize}[leftmargin=1.8em,itemsep=.5em]
 \item \textbf{R 中的公开数据。}在 R 中公开提供的 13 个数据集中，样本量均较小（从 200 到 654），协变量也较少（从 3 到 18），但超级学习器的表现优于所研究的每一个单独算法。这项研究的重要之处在于，它用真实数据表明，超级学习的优势并不以大样本和众多协变量为前提。在所考虑的 20 个单独算法中，参数线性回归的总体表现仅排在第 8 位（Polley and van der Laan，2010；Polley et al.，2011）。
 \item \textbf{死亡风险评分。}这项研究使用超级学习，构造了一个预测老年人群死亡风险的改进函数。它还表明，一项精心收集数据的小型队列研究，尽管其规模仅为大型健康数据库的十分之一（无论按研究对象数还是协变量数衡量），却能得到更准确的预测函数。因此，这篇论文较早揭示了大型电子健康数据库在回答特定科学研究问题时的局限性（Rose，2013）。
 \item \textbf{重症监护病房中的死亡风险评分。}构造重症监护病房（intensive care unit，ICU）的死亡风险评分，是一个难题；以往的评分系统在验证研究中表现不佳。针对这一问题开发的超级学习器表现优异，受试者工作特征曲线（receiver operating characteristic curve）下面积达到 94\%（Pirracchio et al.，2015）。
 \item \textbf{HIV RNA 监测。}这项研究表明，利用电子健康记录中的服药依从性数据实施超级学习，可能有助于识别病毒学失败高风险患者（Petersen et al.，2015）。
 \item \textbf{医保计划支付风险调整。}目前，用于确定向医疗保险计划支付金额的方法，完全是参数化的。近期针对这一问题的超级学习研究表明，利用这一非参数框架选出的简化风险调整公式，保留了传统较大规模公式的大部分效率。这可能影响医疗保险机构通过激进地提高诊断编码等级或实施欺诈来操纵制度的能力（Rose，2016）。
\end{itemize}

超级学习器具有已经确立的理想统计性质，《靶向学习》（2011）第 3 章对此有详细讨论。简言之，在所有候选算法均达不到一个事先正确设定的参数统计模型所具有的收敛速率的现实情形中，超级学习器的渐近表现，与所有可能的加权组合中的最优选择一样好，既包括收敛速率，也包括其中的常数。下面利用 van der Laan and Dudoit（2003）的有限样本预言不等式（finite sample oracle inequality），正式重述这一结论。

\begin{sourcebox}
\noindent\textbf{有限样本预言不等式。}给定一组估计量（即算法）$P_n\mapsto\widehat Q_k(P_n)$，基于损失函数的交叉验证选择器（cross-validation selector）为
\[
 k_n=\widehat K(P_n)=\argmin_k E_{B_n}P^1_{n,B_n}L\bigl(\widehat Q_k(P^0_{n,B_n})\bigr),
\]
其中，$B_n\in\{0,1\}^n$ 是一个随机变量，将数据划分为训练集 $\{i:B_n(i)=0\}$ 和验证集 $\{i:B_n(i)=1\}$；$P^0_{n,B_n}$ 为训练集的经验分布（empirical distribution），$P^1_{n,B_n}$ 为验证集的经验分布，而 $L(\cdot)$ 是我们的差异度量，即损失函数。所得估计量称为离散超级学习器（discrete super learner）：$\widehat Q(P_n)=\widehat Q_{\widehat K(P_n)}(P_n)$。

我们考虑的损失函数满足
\[
 \sup_Q
 \frac{\operatorname{var}_{P_0}\{L(Q)-L(Q_0)\}}
      {P_0\{L(Q)-L(Q_0)\}}
 \leq M_2,
\]
并且一致有界：
\[
 \sup_{O,Q}|L(Q)-L(Q_0)|(O)<M_1<\infty,
\]
其中，上确界遍及 $P_0$ 的支撑集，以及将要考虑的所有可能的 $Q_0$ 估计量。

在损失函数 $L$ 在 $P_0$ 的支撑集上一致有界的假设下，对于固定的 $n$，交叉验证选择器的差异度量与预言选择器（oracle selector）的差异度量之间的余项，在所有数据生成分布上一致成立。这表明，交叉验证选择器能够以 $\log(K(n))/n$ 的距离逼近预言选择器的表现。更确切地说，对于二次损失函数，交叉验证选择器满足以下预言不等式：
% 原书第 31 页左端 L( 的闭括号缺失；按印刷原式保留，不擅自补正。
\[
 \begin{aligned}
 &E_{B_n}\bigl\{P_0L(\widehat Q_{k_n}(P^0_{n,B_n})-L(Q_0)\bigr\}\\
 &\quad\leq(1+2\delta)E_{B_n}\min_k
 P_0\bigl\{L(\widehat Q_k(P^0_{n,B_n}))-L(Q_0)\bigr\}\\
 &\qquad{}+2C(M_1,M_2,\delta)\frac{1+\log K(n)}{np},
 \end{aligned}
\]
其中，$\delta>0$，常数 $C(M_1,M_2,\delta)=2(1+\delta)^2(M_1/3+M_2/3)$，$p$ 是验证集中研究对象所占的比例，$K(n)$ 是算法集合中的算法数。这些结果可以推广到估计得到的损失函数 $L_n(Q)$；这类损失函数逼近某个固定的损失函数 $L(Q)$。
\end{sourcebox}
